\chapter{CPU a architektura počítače}

\textbf{Centrální procesorová jednotka (CPU)} je základní řídicí prvek počítačového systému, jehož
úkolem je vykonávání instrukcí uložených v pamětí a koordinace činnosti ostatních komponentů systému.
Přestože se konkrétní realizace CPU výrazně liší, jejich základní principy a funkce zůstává napříč
architekturami obdobné.

Posloupnost instrukcí určených k vykonání procesorem se nazývá \textbf{strojový kód}
(\textit{machine code}). Tento kód je reprezentován v binární formě a optimalizován pro efektivní
zpracování procesorem, nikoliv pro přímé čtení a úpravy člověkem.

Procesor lze chápat jako \textbf{stavový stroj}, jehož stav je tvořen \emph{registry}, \emph{stavovými
příznaky} a případně dalšími interními prvky. Vykonáním instrukcí dochází ke změně tohoto stavu.
Tento pohled je zejména důležitý při implementaci emulátoru, kde cílem není fyzická simulace
tranzistorů, ale reprodukovat jeho pozorovatelné chování, tedy změnu stavu a interakce s okolními
komponenty systému.

Aby procesor byl více užitečný, je součástí širšího systému, ve kterém může komunikovat s pamětí,
vstupně-výstupnými a jinými druhy zařízení. Prostředek, který spojuje více zařízení a umožňuje
vzájemnou komunikaci, se nazývá \textbf{sběrnice}.

\label{sec:tp-architecture}
\section{Architektura počítače}

Architektura výpočetního systému je obecně definována jako množina dostupných datových typů, operací
a funkcí. Architektura počítačů je pak inženýrským oborem, který se zabývá návrhem a strukturou
těchto systémů. Podle Tanenbauma \cite{tanenbaum} lze moderní počítače rozdělit do několika
hierarchických architektonických úrovní (vrstev):

\begin{enumerate}
    \item \textbf{Číslicové obvody} -- nejnižší úroveň tvořená logickými hradly (AND, OR, NOT).
        Jelikož je cílem práce softwarová emulace zaměřená na replikaci vnějšího chování systému,
        leží tato fyzická vrstva mimo náš přímý zájem.
    \item \textbf{Mikroarchitektura} -- zabývá se realizací funkčních bloků procesoru, jako jsou
        registry a aritmeticko-logická jednotka (ALU). Jejich propojení tvoří \textbf{datovou cestu}
        (\textit{datapath}). Pro emulaci je tato úroveň klíčová z hlediska definice stavu registrů a
        vnitřní paměti.
    \item \textbf{Instrukční sada (ISA)} -- klíčové rozhraní mezi hardwarem a programátorem
        \cite{arm-glossary-isa}. Definuje formát instrukcí, adresní režimy a správu paměti. Precizní
        implementace této vrstvy je pro vývoj emulátoru stěžejní.
    \item \textbf{Operační systém} -- softwarová vrstva spravující prostředky počítače. Architektura
        konzole NES však tradičním operačním systémem nedisponuje a spouští kód přímo „na holém
        železe“ (\textit{bare-metal}), proto tato vrstva v našem případě odpadá.
    \item \textbf{Jazyk symbolických adres (\textit{Assembly})} -- čitelnější reprezentace
        instrukční sady využívající mnemotechnické zkratky (např. \texttt{LDA}, \texttt{STA}). V
        tomto jazyce byla napsána naprostá většina her pro procesor MOS 6502 v systému NES.
    \item \textbf{Vyšší programovací jazyky} -- jazyky jako C/C++ nabízejí vysokou míru abstrakce a
        přenositelnost. Jazyky jako C využívají modernější hry pro NES, a právě na této úrovni je
        vyvíjen náš emulátor pro platformu RP2350.
\end{enumerate}

\section{Organizace CPU}

Drtivá většina architektur jsou založeny na \textbf{von Neumannové architektuře}, která byla popsaná
americkým matematikem Johnem von Neumannem v roce 1945. Tento model je založen na principu
\emph{programovatelného počítače}, popsaný britským matematikem Alanem Turingem v roce 1936, který
umožňoval snadno přeprogramovat počítač, aniž by bylo zapotřebí zasahovat do jeho hardwarové strany.

\begin{figure}[htbp]
    \centering
    \includesvg[width=\textwidth]{images/von_neumann.svg}
    \caption{Von Neumannův model CPU. }
    \label{fig:von_neumann}
\end{figure}

Charakteristickým rysem této architektury je to, že nerozděluje instrukce a data, což ma určité
implikace:

\begin{itemize}
    \item Instrukce jsou reprezentovány stejným způsobem, jako data (binárně)
    \item Paměť pro instrukce a data je společná
    \item Paměť je zpřístupňovaná pomoci adres, nezávisle na typu dat v odpovídající paměťové buňce.
\end{itemize}

Většina architektur využívajících Von Neumannův model se skládají z pěti hlavních funkčních
bloků:

\begin{itemize}
    \item \textbf{Stav CPU}
    \item \textbf{Řídící jednotka (CU)}
    \item \textbf{Aritmeticko-logická jednotka (ALU)}
    \item \textbf{Sběrnice (komunikační rozhraní)}
    \item \textbf{Paměťová jednotka (paměťové rozhraní)}
    \item \textbf{Vstupné zařízení}
    \item \textbf{Výstupné zařízení}
\end{itemize}

\subsection{Řídící jednotka}

Táto jednotka je zodpovědná za celkové řízení procesoru, koordinuje tok dat mezi registry a ostatní
zařízení, a generuje řídící signály. Právě tady je implementován instrukční cyklus, a na základě
aktuální instrukce, určuje jaké operace mají být vykonány v daném kroku.

\subsection{Aritmeticko-logická jednotka}

Aritmeticko-logická jednotka (ALU, anglicky \textit{Arithmetic-Logic Unit}) je zodpovědná za provádění
aritmetických a logických operací. Pracuje s operandy uloženými v registrech CPU, nebo získává je
přímo z pamětí, a výsledek operací přímo ovlivňuje stav procesoru (obsah registru a stavových
příznaků).

V jednoduchých procesorech, jako je např. MOS~6502, představuje ALU hlavní vykonávací jednotku,
která provádí skutečné vykonání většiny instrukcí. I v moderních procesorech zůstává ALU jedním ze
základních stavebních prvků každého jádra, přestože jádro je doplněno o další vykonávací jednotky
(FPU, SIMD apod.). Avšak je to mimo rozsah této práce.

\subsection{Paměťová jednotka}

V Von Neumannové architektuře paměťová jednotka představuje jednotný prostor pro uchovávání dat a
instrukcí, a obvykle je to interní paměť (případ mikrokontrolérů), nebo externí (případ osobních
počítačů).

Pro adresování pamětí, CPU obsahuje paměťové rozhraní, které zajišťuje přístup k pamětí a paměťové
mapovaným zařízením (viz sekce \ref{sec:tp-bus} o sběrnici). Umožňuje čtení dat a zápis výsledků
výpočtu zpět do pamětí. V jednoduchých architekturách může jít pouze o logiku adresování, zatímco v
komplexnějších architekturách paměťové rozhraní může zahrnovat i složitější správu pamětí.

\label{sec:tp-bus}
\subsection{Komunikační rozhraní}

Komunikační rozhraní, neboli \textbf{sběrnice} (anglicky \textit{bus}), slouží jako základní
mechanizmus pro propojení více zařízení a zajištění přenášení dat, adres a řídicích signálů mezi
procesorem, pamětí a periferiemi. Sběrnici lze chápat jako fyzické i logické propojení jednotlivých
komponentů systému, a spolu s paměťovým rozhraním, tvoří \textbf{most} mezi procesorem a vnějším
světem.

\subsection{Von Neumannovo úzké hrdlo (bottleneck)}

Ačkoliv sjednocení pamětí pro instrukce a data přináší flexibilitu, nese s sebou i hlavní nevýhodu
tohoto modelu, známou jako von Neumannovo úzké hrdlo.

Instrukce a data se nacházejí ve stejné pamětí, musí sdílet společnou komunikační cestu (sběrnici)
spojující paměť s procesorem. V praxi to znamená, že pokud procesor v daný okamžik zapisuje nebo čte
data z pamětí, nemůže paralelně po stejné sběrnici načítat další instrukci. Procesor tak musí čekat
na dokončení paměťových transakcí, což omezuje jeho potenciální rychlost. Tento problém řeší
alternativní modely CPU, jako \emph{Harvardská architektura} nebo hybridní modely.

\section{Alternativní modely CPU}

Výše popsáná von Neumannová architektura není jediný způsob organizace CPU.

\subsection{Harvardská architektura}

Harvardská architektura striktně odděluje paměť na dvě fyzicky i logicky nezávislé části: jednu pro
instrukce a druhou pro data. Každá paměť má svou vlastní adresovou i datovou sběrnici. Díky tomu
může procesor v jednom taktu současně načítat novou instrukci z programové paměti a zároveň číst
nebo zapisovat data do paměti datové, což výrazně zvyšuje výkon systému.

Existence Harvardské architektury jako samostatné koncepce byla některými výzkumníky
zpochybněna\cite{harvard-myth}. Vzhledem k tomu, že tento koncept je poměrně rozšířený, tato práce
se rozhodla jej popsat jako samostatný model CPU.

\subsection{Modifikované a hybridní modely}

V současné počítačové architektuře se nejčastěji setkáváme s kombinací obou výše zmíněných přístupů,
která se označuje jako \textbf{Modifikovaná Harvardská architektura}. Tento model se snaží
eliminovat nevýhody obou koncepcí: zachovává vysoký výkon Harvardu, ale zároveň nabízí flexibilitu
von Neumannovy architektury.

Hlavním rysem hybridních modelů je existence oddělených vyrovnávacích pamětí (Cache) na nejnižší
úrovni procesoru (L1 cache). Procesor má vnitřně separátní instrukční cache a datovou cache, což mu
umožňuje přistupovat k oběma typům informací současně a efektivně využívat proudové zpracování
instrukcí (\textit{pipeline}).

Směrem k hlavní operační paměti se však tyto dvě cesty sbíhají do jedné. Z pohledu programátora se
tedy paměť jeví jako jednotný adresní prostor, kde jsou instrukce i data uloženy společně. Mezi
hlavní výhody tohoto hybridního přístupu patří:

\begin{itemize}
    \item Díky odděleným cache nedochází k „úzkému hrdlu“ při současném čtení instrukcí a manipulaci
        s daty.
    \item Operační systém může snadno měnit rozdělení paměti mezi program a data, což by u čistého
        Harvardu vyžadovalo fyzickou změnu hardwaru.
    \item Program může zapisovat data do paměti, která jsou následně interpretována jako instrukce
        (např. při JIT kompilaci).
\end{itemize}

Typickým příkladem této architektury jsou moderní procesory rodiny x86 nebo jádra ARM Cortex, na
kterých je založen i mikrokontrolér RP2350 použitý v tomto projektu. Ačkoliv se navenek tváří jako
systémy s jednotnou pamětí, jejich vnitřní organizace využívá výhod Harvardského modelu pro dosažení
maximálního výpočetního výkonu.

\section{Stav CPU}

Stav CPU představuje úplný popis vnitřního stavu procesoru v daném okamžiku. Je tvořen
\textbf{registry}
a \textbf{stavovými příznaky}. Vykonáním instrukce dochází ke změně tohoto stavu.

Z pohledu emulace je stav CPU klíčovým k tomu, aby správně napodobit chování procesoru. Cílem
přesného emulátoru je reprodukovat tyto změny stavu v reakci na vykonávané instrukce a vstupy z
okolních zařízení.

\subsection{Registry}

Procesorový registr je vysokorychlostní paměť, která je součástí procesoru a je vždy mu přístupná.

Registry nemusí být jenom pamětí navíc, ale také mohou mít i speciální vlastností nebo funkce,
nebo se lišit podle dovoleného způsobu přístupu:

\begin{itemize}
    \item Uživatelsky přístupné
        \begin{itemize}
        \item \textbf{Datové registry} uchovávají celočíselné hodnoty, v pokročilých architekturách
            i čísla s plovoucí tečkou, a také znaky nebo malé bitové pole.
            \begin{itemize}
                \item Do této skupiny také patří registry zásobníku, jako je vrchol zásobníku nebo
                    registr uchovávající jeho kopii (viz kapitola \ref{sec:tp-stack}).
            \end{itemize}
        \item \textbf{Adresové registry} uchovávají adresy pro instrukce, které přistupují paměť.
        \item \textbf{Univerzální (general-purpose) registry} mohou uchovávat i data, i adresy.
        \item \textbf{Stavové (status) registry} vyznačují příznaky různých stavů procesoru.
        \item \textbf{Speciální (special-purpose) registry}
            \begin{itemize}
                \item \textbf{Registr čítače instrukcí PC} obsahuje adresu aktuální instrukce (nebo
                    další instrukce k vykonání)
                \item \textbf{Registry tykající se zásobníku} taky se někdy zahrnují do této skupiny.
            \end{itemize}
        \item \textbf{Řídící (control) registry} zpravidla obsahují nastavení procesoru, ovlivňující
            jeho funkčnost a proces vykonávání instrukcí (např., jeho taktová rychlost, přerušení
            apod.).
        \end{itemize}
    \item Interní registry
        \begin{itemize}
            \item \textbf{Instrukční registr (IR)} obsahuje kód aktuálně vykonávané instrukce.
            \item Registry, určené k dočasnému uchovávání různých informací.
        \end{itemize}
\end{itemize}

Interní registry jako čítač instrukcí PC, instrukční registr IR pro uchovávání kódu instrukcí, jsou
obsažené v každém procesoru pro jeho minimální funkčnost.

\label{sec:tp-stack}
\subsection{Zásobník}

\begin{figure}[htbp]
    \centering
    \includesvg[width=0.75\textwidth]{images/stack.svg}
    \caption{Znázornění zásobníku v pamětí, kde typicky je mapován.}
    \label{fig:stack}
\end{figure}

Zásobník (\textit{stack}) je buď samostatná paměť, nebo vyhrazená oblast v interní pamětí,
která slouží k dočasnému uchovávání informací (např. stavový registr před vykonáním instrukce,
hodnota PC před přerušením nebo změnou PC).

Zásobník je realizován jako obecná datová struktura, pro kterou je charakteristický takový způsob
manipulace s data, že data uložena jako poslední budou čtena jako první (\textit{LIFO - Last In,
First Out}). Pro manipulaci s uloženými datovými položkami se udržuje tzv. \textbf{ukazatel
zásobníku}, neboli \textbf{vrchol zásobníku}. Ten odkazuje na poslední položku. Obdobně procesor
typicky také udržuje registr, který obsahuje adresu vrcholu zásobníku.

Přítomnost zásobníku je klíčové pro rozsáhlé programy, využívající větvení. Navíc zásobník může
sloužit jako paměť navíc, např. pro uchovávání více mezivýsledků nebo většího počtu vstupních dat.

\section{Instrukce a instrukční sada}

Jak již bylo zmíněno v sekci \ref{sec:tp-architecture}, instrukční sada definuje programovací
rozhraní mezi programátorem a hardwarem procesoru. Instrukční sada typicky zahrnuje tyto komponenty:

\begin{itemize}
    \item Sadu dostupných instrukcí: seznam všech instrukcí podporovaných procesorem.
    \item Format instrukcí: rozložení na úrovni bitů definující operační kódy, operandy a adresovací
        režimy.
    \item Registry: jejích počet, druhy, a jejích účely.
    \item Datové typy: druhy dat, se kterými procesor může přímo pracovat (celá čísla, plovoucí
        čísla apod.).
    \item Architektura paměti: endianita více bajtových dat, mapování, virtuální paměť, ochrana
        určitých oblastí v pamětí apod.
    \item Přerušení a výjimky.
\end{itemize}

Instrukce slouží k řízení CPU, a lze je chápat jako změnu stavu; vykonáním instrukce dochází ke
změně stavu jednoho nebo více registrů, stavových příznaků nebo dochází k interakci s pamětí a
jinými zařízeními. Výsledný stav je určen předchozím stavem CPU a vykonanou instrukci.

Instrukce jsou typicky uloženy v externí pamětí (někdy v interní, např. rodina 8-bitových
mikrořadičů AVR), odkud jsou postupně čteny a vykonávány.

\subsection{Data a operandy}

Většina instrukcí - zejména aritmetických, logických, nebo paměťových - potřebuji vstupní data k
vykonání své činností. Data, která jsou součástí zakódované instrukce, se nazývají \textbf{operandy}.
Druhy vstupních dat neboli operandů lze rozdělit do 3 základních typů, které určují typ dat, a
také jejích směr nebo původ:

\begin{itemize}
    \item Konstanta (číslo)
    \item Adresa do pamětí
    \item Číslo nebo název registru
\end{itemize}

Podle směru a původu vstupních dat, lze instrukce rozdělit do následujících kategorií:

\begin{itemize}
    \item \textbf{Read} instrukce
    \item \textbf{Write} instrukce
    \item \textbf{Read-modify-write} instrukce
\end{itemize}

\subsection{Adresovací režimy}

S přihlédnutím k tomu, že je několik možných zdrojů vstupních dat (konstanta, číslo registru nebo
paměťová adresa), je několik způsobů tyto data získat v průběhu dekódování instrukce. Způsob,
kterým procesor získá vstupní data při dekódování instrukce, se nazývá \textbf{adresovací režim}.

Určitý počet adresovacích režimů a jejích význam se značně liší u různých architektur. Jsou
architektury (např. většina RISC architektur), které preferují co nejmenší počet adresovacích
režimů, zatímco jsou architektury (CISC architektury jako DEC VAX) které mohou mít je více než
10.

Dále budou popsány nejrozšířenější adresovací režimy, které se vyskytují ve většině existujících
architektur:\cite{tanenbaum}

\begin{itemize}
    \item \textbf{Okamžité adresování (immediate addressing)} - nejjednodušší způsob určit operand
        je uchovávat ho jako součást instrukce, než adresu operandu nebo jinou informaci, určující
        jeho umístění. Hlavní nevýhodou je omezení jenom na konstantní hodnoty.
    \item \textbf{Přímé adresování (direct addressing)}, nebo také \textbf{absolutní adresování
        (absolute addressing)} - adresa operandu je uchovávaná spolu s instrukcí. Při dekódování,
        procesor přečte přímo tuto adresu aby získal operand.
    \item \textbf{Registrové adresování (register addressing)} - konceptuálně podobné jako
        přímé adresování, ale místo adresy do pamětí se specifikuje číslo registru.
    \item \textbf{Registrové nepřímé adresování (register indirect addressing)} - tento režim
        je sjednocení přímého a registrového adresování: místo přímého uchovávání výsledné adresy
        operandu spolu s instrukcí, adresa operandu je uchovávána ve specifikovaném registru.
    \item \textbf{Indexované adresování (indexed addressing)} - adresa operandu je uchovávaná v
        registru, a při její načtení, k výsledné adrese se přičte konstantní celočíselné číslo
        (posun).
    \item \textbf{Indexované adresování pomocí základny s posunem (based-indexed addressing)} -
        podobně jako indexované adresování, tady výsledná adresa operandu se získá přičtením
        základny z registru a konstantního posunu. V tomto adresování ale posun je uchováván také
        v registru, tím pádem toto adresování v některých případech může být flexibilnější, když
        posun není předem známý.
\end{itemize}

\label{subsec:tp-instruction-cycle}
\subsection{Instrukční cyklus}

Po inicializaci, CPU se nachází ve stavu, který se nazývá instrukční cyklus (anglicky
\emph{fetch-decode-execute cycle}). Je to neustálý proces, když procesor načítá další instrukcí
(fetch), zpracovává jí (decode) a pak vykonává (execute). Instrukční cyklus se vykonává sekvenčně,
každá instrukce je zpracována před spuštěním další instrukce.

Instrukční cyklus lze chápat jako opakovanou posloupnost kroků, při nichž procesor na základě svého
aktuálního stavu načte další instrukci, určí její význam a vykoná instrukci - provede odpovídající
změny vnitřního stavu nebo interakce s okolními komponentami systému.
Z pohledu emulace je důležité, že tento proces je deterministický: při stejném vstupním stavu a
stejném obsahu paměti a případně při stejném stavu okolních komponentů, musí procesor vždy
vykazovat shodné chování.

Ačkoli konkrétní implementace instrukčního cyklu se mezi architekturami liší (zejména u moderních
procesorů, kde instrukční cyklus může se provádět paralelně; viz \emph{instruction pipeline}), lze
obecněji instrukční cyklus rozdělit do následujících fází:

\begin{enumerate}
    \item \textbf{Načtení instrukce (fetch)}: tato fáze zahajuje instrukční cyklus získáním kódu
        instrukce z pamětí. Podle aktuální hodnoty čítače instrukcí PC, procesor se zeptá na adres
        následující instrukce v pamětí, pak je instrukce uložena do interního registru IR. PC je
        aktualizován tak, aby ukazoval na další instrukcí.
    \item \textbf{Dekódování instrukce (decode)}: během této fáze, uložena instrukce v IR je
        analyzovaná řídící jednotkou. Procesor urči typ instrukce (aritmetika, logika, paměťová
        operace), adresovací režim, a vstupní data (operandy) zapotřebí k vykonání této instrukce.
    \item \textbf{Vykonání instrukce (execute)}: vykonání představuje vlastní realizaci významu
        instrukce. Může zahrnovat aritmetické a logické operace, přístup do pamětí, změna stavu
        registrů, úpravu stavových příznaků nebo komunikaci s periferiemi. Je důležité ale
        poznamenat, že některé instrukce se vykonávají v jednom kroku, některé ve více krocích.
\end{enumerate}

Z hlediska programátora, instrukční cyklus se vykonává sekvenčně: další cyklus se začne až po
dokončení předchozího cyklu.

\subsection{Třídy instrukcí}

Podle činnosti, původu dat a směru výstupních dat (výsledku), instrukce se obecně rozděluji na:

\begin{itemize}
    \item \textbf{Datové a paměťové instrukce}
        \begin{itemize}
            \item Nastavení hodnoty registru na konstantní hodnotu.
            \item Kopírování dat z jednoho místa (registru nebo pamětí) do jiného.
            \item Operace se zásobníkem, např. Přidání datové položky, odebraní nebo získaní aktuálních
                dat v místě, kam ukazuje vrchol zásobníku.
        \end{itemize}
    \item \textbf{Aritmetické a logické instrukce}
        \begin{itemize}
            \item Základní aritmetické operace jako sčítání, odčítání, násobení, nebo dělení. Typicky,
                tyto instrukce nastavují stavové příznaky podle hodnoty výsledku nebo jiných událostí.
            \item Bitové operace jako je konjunkce (AND), disjunkce (OR), znegování (NOT).
            \item Porovnání obsahu dvou registrů.
        \end{itemize}
    \item \textbf{Řídící instrukce}
        \begin{itemize}
            \item Větvení programu (změna registru čítače instrukcí PC).
            \item Podmínečné větvení programu nebo přeskakování další instrukce (např. podle aktuálních
                stavových příznaků).
        \end{itemize}
\end{itemize}

U složitých architektur rozdělení instrukcí může být rozsáhlejší.

\section{Časování a synchronizace}

Podobně jako většina digitálních obvodů, procesor je řízen časem, který je určen periodickým
hodinovým signálem, označovaným jako \textbf{takt} nebo \textbf{hodiny} (\textit{clock}) neboli
cyklus (obr. \ref{fig:cpu_clock}). Taktování určuje, v jaký čas mohou jednotlivé částí procesoru
měnit svůj stav. Obecně takt ovlivňuje rychlost (přesněji frekvenci) vykonávání instrukcí: čím větší
takt, tím víc instrukcí se vykoná za jednotku času.

\begin{figure}[htbp]
    \centering
    \includesvg[width=\textwidth]{images/cpu_clock.svg}
    \caption{Znázornění taktování CPU. V tomto případě se jedna o instrukcí, která zapisuje výsledek
    zpět do paměti.}
    \label{fig:cpu_clock}
\end{figure}

Takt rozděluje čas na diskrétní kroky, ve kterých dochází k synchronizované změně stavu procesoru,
tj. registrů a stavových příznaků. Takt zajišťuje, že všechny jednotky procesoru jsou ve stejném
časovém rámci, tím pádem výsledky operací jsou konzistentní a předvidatelné.

Procesor, který vykonává instrukce ve více cyklech, se nazývá multi-cycle procesor: každý cyklus
odpovídá dílčí částí instrukce, např. přístup do pamětí. nebo zpracování mezivýsledků. Multi-cyklový
model procesoru má výhody oproti jednocyklovému modelu: komplexnější instrukce zabírají dost
procesorového času, ale běžné instrukce zabírají menší čas, a také možnost opětovného a efektivního
využití existujících jednotek.

Klíčovým efektem synchronizace je determinismus, neboli předvidatelnost procesoru a celého systému
jako celku. Předvidatelnost znamená, že při stejném vstupu a počátečním stavu výsledek bude vždy
časově a funkčně stejný. Zejména v systémech, které obsahují více procesoru, grafický nebo zvukový
subsystem apod., je přesné časování zásadní, protože činnost každého komponentu v podobném systému
probíhá paralelně. Jakákoli odchylka v časování může vést k nepřesnosti a nečekaným výsledkům.

\section{Přerušení}

% \begin{figure}[H]
%     \centering
%     \includesvg[width=\textwidth]{images/cpu_int.svg}
%     \caption{Integrace přerušení do instrukčního cyklu.}
%     \label{fig:cpu_int}
% \end{figure}

Přerušení představují mechanismus, kterým může být běžný sekvenční tok vykonávání instrukcí dočasně
pozastaven v reakci na vnější událost (stisknutí tlačítka, změna napěťového úrovně apod.) nebo
vnitřní (závažná chyba za běhu, žádost o posílaní nových dat apod.).

Díky tomuto mechanizmu, procesor nemusí obětovat svůj čas na dotazování všech možných zdrojů
událostí, a navíc tento model má nevýhodu, že např. hned jak procesor ukončil dotazování, může
nastat vnější událost, která ale skončí před opětovným dotazováním, a procesor nebude schopen tuto
událost zaznamenat.

Aby program mohl přerušení zpracovat, musí mít přiřazený tzv. \textbf{obslužnou rutinu}
(\textit{interrupt handler} nebo \textit{interrupt service routine}, zkráceně ISR), což je kód,
který je vyvolán v okamžik vzniku tohoto přerušení. Když přerušení nastává, současný stav procesoru
je uložen (zpravidla do zásobníku), včetně čítače instrukcí PC a dalších stavových registrů, a
procesor začíná vykonávat obslužnou rutinu, přiuzenou přerušení. Pokud přerušení nastane uprostřed
instrukčního cyklu, začne se zpracovávat \textbf{až po vykonávání aktuální instrukce}. Vnitřně
procesor eviduje požadavky ne přerušení, a před zahájením dalšího instrukčního cyklu, rozhodne, zda
přerušení má být zpracováno. Po dokončení obsluhy přerušení, procesor se vrátí k původnímu programu
a~obnoví uložený stav a~pokračuje v instrukčním toku.

Rozdělení přerušení:

\begin{itemize}
    \item \textbf{Vnější přerušení} jsou generovány z vně, např. periferiemi.
    \item \textbf{Vnitřní přerušení} jsou vyvolány samotným programem během vykonávání.
    \item \textbf{Maskovatelná přerušení}, která lze dočasně zakázat (ignorovat) nebo povolit.
    \item \textbf{Nemaskovatelná přerušení}, která mají vždy přednost před ostatními přerušení,
        a nezle je ignorovat.
\end{itemize}

% Tento mechanizmus umožňuje asynchronní programování a je podobný událostmi řízenému
% programování (\textit{event-driven programming}), tj. styl programování, kde podobně každá událost
% je zpracovaná obslužní funkcí.
%
% Avšak mezi přerušení a tímto stylem programování existují zásadní rozdíly:
%
% \begin{itemize}
%     \item
%     \item přerušení přeruší vykonávání v aktuálním bodě programu, a obslužní rutina přerušení se
%         začne vykonávat, zatímco event-driven programování má stanovený čas kdy se událostí zpracují.
% \end{itemize}
